\begin{textsample}{POS Dim 3 – ac – Score 53.00 – ac/ac\_inf\_23.txt}  \label{ex:f3_pos_021}
23. \textbf{CRIANÇAS} \textbf{PEQUENAS} ( 4 ANOS A 5 ANOS E 11 \textbf{MESES} ) 23.1.
Campo de experiências “ Corpo, \textbf{gestos} e movimentos ” Com o corpo ( por meio dos sentidos, \textbf{gestos}, movimentos impulsivos ou intencionais, coordenados ou espontâneos ), as \textbf{crianças}, desde cedo, exploram o mundo, o espaço e os objetos do seu entorno, estabelecem relações, expressam -se, \textbf{brincam} e produzem conhecimentos sobre si, sobre o outro, sobre o universo social e cultural, tornando -se, \textbf{progressivamente}, conscientes dessa corporeidade.
Por meio das diferentes linguagens, como a música, a dança, o teatro, as \textbf{brincadeiras} de faz de conta, elas se comunicam e se expressam no entrelaçamento entre corpo, emoção e linguagem.
As \textbf{crianças} conhecem e reconhecem as sensações e funções de seu corpo e, com seus \textbf{gestos} e movimentos, identificam suas potencialidades e seus limites, desenvolvendo, ao mesmo tempo, a consciência sobre o que é seguro e o que pode ser um risco à sua integridade física.
Na Educação \textbf{Infantil}, o corpo das \textbf{crianças} ganha centralidade, pois ele é o partícipe privilegiado das práticas pedagógicas de \textbf{cuidado} físico, orientadas para a emancipação e a liberdade, e não para a submissão.
Assim, a instituição escolar precisa promover oportunidades ricas para que as \textbf{crianças} possam, sempre animadas pelo espírito \textbf{lúdico} e na interação com seus pares, explorar e vivenciar um amplo repertório de movimentos, \textbf{gestos}, olhares, \textbf{sons} e \textbf{mímicas} com o corpo, para descobrir variados modos de ocupação e uso do espaço com o corpo ( tais como sentar com apoio, rastejar, engatinhar, escorregar, caminhar \textbf{apoiando} -se em berços, mesas e cordas, saltar, escalar, equilibrar -se, correr, dar cambalhotas, alongar -se etc. ). 23.2.
Direitos de aprendizagem e desenvolvimento - Conviver com \textbf{crianças} e \textbf{adultos}, experimentando as possibilidades corporais para interagir e desenvolver o interesse sobre os \textbf{cuidados} com o corpo. - \textbf{Brincar} utilizando \textbf{criativamente} o repertório da cultura corporal e do movimento, expressando emoções, alegrias e (re)significando o mundo a sua \textbf{volta}. - Explorar os movimentos corporais através de \textbf{brincadeiras} e interações para expressar suas emoções e desenvolver sua autonomia. - Participar de atividades que envolvem práticas corporais como dança, música, teatro, artes circenses, manifestando encanto e satisfação, desenvolvendo a autonomia. - Expressar corporalmente emoções e representações tanto nas relações cotidianas como nas \textbf{brincadeiras}, artes visuais, dramatizações, danças, músicas, \textbf{escuta} e contação de histórias. - Conhecer -se por meio de diversas experiências de interações, \textbf{brincadeiras} e explorações com seu corpo, desenvolvendo a autonomia e a autoestima. 23.2.1.
Objetivo de aprendizagem e desenvolvimento : ( EI03CG01 ) - Criar com o corpo formas diversificadas de expressão de sentimentos, sensações e emoções, tanto nas situações do cotidiano quanto em \textbf{brincadeiras}, dança, teatro, música.
Aprendizagens específicas - Expressar -se intencionalmente de diferentes modos através de \textbf{gestos}, expressões corporais e movimentos do corpo. - Manifestar afetividade em diferentes situações do cotidiano. - Conhecer o próprio corpo, aceitando e valorizando suas características corporais, expressando -se de diferentes formas e construindo uma imagem positiva de si. - Expressar diferentes percepções, sensações e emoções, proporcionadas pelos órgãos dos sentidos, estimulando a memória visual, auditiva, olfativa, tátil e gustativa. - Expressar sentimentos, sensações e emoções, por meio das diferentes práticas e expressões corporais. - Perceber e distinguir emoções e sentimentos, em si mesma e nos outros, nas diferentes situações do cotidiano.
Sugestões de experiências - Vivenciar diferentes práticas da cultura corporal e compartilhar sentimentos e as dificuldades encontradas. - \textbf{Brincar} em frente ao espelho, percebendo as diferenças do próprio corpo e do outro. - Criar movimentos, \textbf{gestos}, olhares, \textbf{mímicas} e \textbf{sons} com o corpo, em \textbf{brincadeiras}, jogos e atividades artísticas. - Fantasiar -se na frente do espelho para \textbf{brincar} com os \textbf{colegas}. - \textbf{Brincar} de \textbf{esconder} e encontrar objetos em diferentes locais, seguindo a descrição da posição destes ( longe, perto, à frente, atrás, ao lado, esquerda, direita ). - \textbf{Brincar} com jogos de percurso, elaborando de trajetos com pontos de partida e chegada. - \textbf{Brincar} com jogos que favoreçam a expressividade corporal. - Dançar, \textbf{dramatizar} e \textbf{brincar}, expressando \textbf{desejos}, sentimentos e ideias. - Vivenciar práticas da cultura popular : festa junina, capoeira, \textbf{brinquedos} cantados, dentre outras. 23.2.2.
Objetivo de aprendizagem e desenvolvimento : ( EI03CG02 ) - \textbf{Demonstrar} controle e adequação do uso de seu corpo em \textbf{brincadeiras} e jogos, \textbf{escuta} e reconto de histórias, atividades artísticas, entre outras possibilidades.
Aprendizagens específicas - Ampliar os limites e as potencialidades do seu corpo, explorando diferentes movimentos físicos e motores, como força, velocidade, resistência e flexibilidade. - Adequar seus movimentos corporais aos de seus \textbf{colegas} em situações de \textbf{brincadeiras}, jogos ou atividades coletivas. - Desenvolver, \textbf{progressivamente}, o controle do corpo, movimentando -se seguindo a orientação do professor, de outras \textbf{crianças} ou criando suas próprias orientações. - Desenvolver \textbf{equilíbrio} estático e dinâmico, explorando diferentes posturas corporais. - Desenvolver o gosto e interesse por apreciar e participar de jogos, \textbf{brincadeiras}, \textbf{escuta} e reconto de histórias e atividades artísticas.
Sugestões experiências - \textbf{Brincar} de elaborar circuitos com diferentes materiais. - \textbf{Brincar} utilizando o corpo, percebendo suas potencialidades e limites. - Construir jogos utilizando suas habilidades manuais no atendimento adequado a seus interesses e necessidades. - Planejar a disposição, organização e reorganização dos materiais. - \textbf{Montar} e desmontar atividades diversas : cenários, jogos, dentre outras. - Jogar e \textbf{brincar} representando personagens no faz de conta e no reconto de histórias. - \textbf{Dramatizar} diversas versões de histórias \textbf{infantis}. - Cantar e recriar cantigas, \textbf{batendo} \textbf{palmas}, assobiando e outras possibilidades. - Andar como robôs, zumbis, gatinhos, maria-mole, dentre outras formas. - \textbf{Brincar} de esconde-esconde, jogar bolar, pega-pega, de seguir o mestre, estátua, caça ao tesouro, cabra-cega, pião, amarelinha e outras. 23.2.3.
Objetivo de aprendizagem e desenvolvimento : ( EI03CG03 ) - Criar movimentos, \textbf{gestos}, olhares e \textbf{mímicas} em \textbf{brincadeiras}, jogos e atividades artísticas como dança, teatro e música.
Aprendizagens específicas - Reconhecer o corpo como instrumento expressivo, capaz de criar movimentos, \textbf{gestos} e construir novos conhecimentos. - Ampliar o repertório de manifestações culturais, estéticas e artísticas, relacionadas aos movimentos corporais. - Desenvolver a capacidade de imaginar, criar e se expressar corporalmente por meio de \textbf{gestos}, \textbf{mímicas}, teatro e dança. - Conhecer diferentes espaços temporais a fim de favorecer a realização de movimentos no que diz respeito à duração, sucessão dos acontecimentos, pausa, velocidade e estrutura rítmica. - Descobrir as possibilidades motoras e expressivas de seu corpo, apropriando -se de jogos e \textbf{brincadeiras} folclóricas, tradicionais e contemporâneas. - Desenvolver o crescente domínio dos movimentos corporais de diferentes formas e em diversos espaços. - Combinar seus movimentos corporais com os de outras \textbf{crianças}, criando novos movimentos, usando \textbf{gestos}, movimentos do seu corpo e sua voz.
Sugestões de experiências - Participar de jogos, danças, \textbf{brincadeiras}. - Cantar e dançar, individualmente ou em grupo, acompanhando o ritmo da música ou criando novos ritmos, utilizando \textbf{mímicas}, \textbf{gestos} e diversos movimentos corporais. - \textbf{Brincar} com a música, \textbf{imitando}, \textbf{inventando}, criando e reproduzindo criações musicais. - \textbf{Brincar} de \textbf{bater} \textbf{palmas}, \textbf{pés}, mãos na barriga, \textbf{dedos} na bochecha, estalar os \textbf{dedos}, dentre outros. - \textbf{Brincar} de capoeira, dançar Carimbó, dentre outros. - \textbf{Brincar} de circo, \textbf{imitando} palhaços, malabaristas, equilibristas, mágicos, dentre outros. - Explorar o espaço de diferentes formas : pular, saltar, dançar combinando ou criando movimentos. 23.2.4.
Objetivo de aprendizagem e desenvolvimento : ( EI03CG04 ) - Adotar \textbf{hábitos} de autocuidado relacionados a \textbf{higiene}, alimentação, \textbf{conforto} e aparência.
Aprendizagens específicas - Reconhecer o próprio corpo, suas partes e segmentos, \textbf{demonstrando} cada vez mais \textbf{interesse} em realizar de forma autônoma ações de autocuidado. - Apropriar -se de procedimentos relativos ao autocuidado, valorizando atitudes relacionadas ao bem-estar, à saúde, à \textbf{higiene}, à alimentação, ao \textbf{conforto}, à segurança e à proteção do corpo. - Compreender os limites corporais para proteger a sua integridade física, desenvolvendo autonomia de movimentos necessários ao autocuidado. - Desenvolver interesse em participar da \textbf{higiene} e organização dos espaços coletivos da Instituição Educativa e de outros espaços, para seu bem-estar e de seu grupo de convivência. - Apropriar -se de \textbf{hábitos} de \textbf{higiene} para promoção da saúde e prevenção de doenças.
Sugestões de experiências - \textbf{Conversar} sobre \textbf{cuidados} com o corpo e a saúde, prevenção de acidentes. - Lavar as mãos antes das \textbf{refeições} e após o uso do banheiro, evitando o \textbf{desperdício} de água. - Cuidar dos dentes após as \textbf{refeições}, pentear os cabelos e outras situações de \textbf{higiene} corporal. - Usar com autonomia diferentes objetos : tesoura, cola, pincel, papéis etc. - Zelar pela própria segurança e dos \textbf{colegas}. - Usar com autonomia diversos materiais de \textbf{higiene} corporal, evitando o \textbf{desperdício}. - Alimentar -se \textbf{autonomamente}, escolhendo os alimentos de sua preferência, servindo -se e evitando o \textbf{desperdício}. 23.2.5.
Objetivo de aprendizagem e desenvolvimento : ( EI03CG05 ) - Coordenar suas habilidades manuais no atendimento adequado a seus interesses e necessidades em situações diversas Aprendizagens específicas - Desenvolver e utilizar habilidades motoras para \textbf{manusear} diferentes objetos e materiais. - Ampliar, \textbf{progressivamente}, o controle de suas habilidades manuais, explorando objetos e materiais de diferentes características ( formas, \textbf{pesos}, \textbf{texturas}, tamanhos e outros ). - Ampliar suas possibilidades de \textbf{manuseio} de diferentes materiais e objetos que favoreçam os movimentos de preensão, \textbf{encaixe} e lançamento. - Desenvolver, gradativamente, a destreza, \textbf{manipulando} \textbf{pequenos} objetos, construindo \textbf{brinquedos} ou jogos e utilizando materiais convencionais e \textbf{reutilizáveis}. - Manifestar sua autonomia ao construir representações gráficas e \textbf{plásticas} ( desenho, \textbf{pintura}, colagem, dobradura, escultura ), explorando materiais diversos. - Perceber os avanços e conquistas relacionadas às suas habilidades manuais.
Sugestões de experiências - \textbf{Brincar} de \textbf{montar}, desmontar, \textbf{empilhar}, encaixar e equilibrar. - Pegar, lançar, amassar, pintar, \textbf{recortar}, rasgar e modelar diversos materiais. - \textbf{Brincar} em atividades que envolvam força, \textbf{equilíbrio} e agilidade. - Criar \textbf{gestos} e movimentos com as mãos e demais membros do corpo. - \textbf{Manipular} objetos, \textbf{brinquedos}, bonecos, fantoches e dedoches em jogos teatrais. 23.3.
Avaliação - \textbf{Acompanhamento} do processo de aprendizagem e desenvolvimento das \textbf{crianças} \textbf{pequenas} : Observação \textbf{criteriosa}, \textbf{registro} e análise quanto : - À utilização adequada das variações do movimento, como : velocidade, lateralidade, força e outros. - À participação dos circuitos explorando movimentos cada vez mais complexos. - À autonomia com relação à destreza e \textbf{equilíbrio} do movimento na hora da \textbf{brincadeira}, da dança, dos jogos, dentre outras. - Às práticas construídas com relação ao autocuidado, saúde, \textbf{higiene} e alimentação. - À atitude de reconhecimento de seus próprios limites e potencialidades ao participar das atividades.

% matched lemmas: acompanhamento, adulto, apoiar, autonomamente, bater, brincadeira, brincar, brinquedo, colega, conforto, conversar, criança, criativamente, criterioso, cuidado, dedo, demonstrar, desejo, desperdício, dramatizar, empilhar, encaixe, equilíbrio, esconder, escuta, gesto, higiene, hábito, imitar, infantil, interessar, inventar, lúdico, manipular, manusear, manuseio, montar, mês, mímico, palma, pequeno, peso, pintura, plástico, progressivamente, pé, recortar, refeição, registro, reutilizável, som, textura, volta
\end{textsample}
